\documentclass[12pt]{article}
\usepackage[T1]{fontenc}
\usepackage[utf8]{inputenc}
\usepackage{lmodern}
\usepackage{textcomp}
\usepackage{enumitem}
\usepackage{geometry}
\usepackage{graphicx}
\usepackage{amsmath, amssymb}
\geometry{margin=1in}
\begin{document}
\begin{center}
Philosophy - Graduate Level Assessment 19 \
Total Marks: 40 \
Answer all questions.
\end{center}

\section*{Multiple Choice (10 marks)}
\begin{enumerate}[label=\arabic*.]
\item Anscombe claims that the notion of moral obligation is derived from the concept of:
\begin{enumerate}[label=(\alph*)]
    \item preference.
    \item natural law.
    \item self-interest.
    \item maximizing utility.
    \item ethical relativism.
\end{enumerate}
\item Berkeley insists that heat and cold are \_\_\_\_\_.
\begin{enumerate}[label=(\alph*)]
    \item elements of nature that do not exist independently
    \item only things existing apart from our minds
    \item only sensations existing in our minds
    \item physical objects
    \item manifestations of our subconscious
\end{enumerate}
\item Aesthetics deals with objects that are\_\_\_\_\_.
\begin{enumerate}[label=(\alph*)]
    \item not essential to our existence
    \item frequently used in daily life
    \item not visible to the human eye
    \item only appreciated by experts
    \item universally liked
\end{enumerate}
\item According to Stevenson, the word "good" has a pleasing emotive meaning that fits it for:
\begin{enumerate}[label=(\alph*)]
    \item dynamic use.
    \item descriptive use.
    \item inferential use.
    \item propositional use.
    \item expressive use.
\end{enumerate}
\item Nathanson proposes a form of retributivism that requires all of the following except
\begin{enumerate}[label=(\alph*)]
    \item a list of crimes ranked according to their seriousness.
    \item a scale of punishments that correspond to the seriousness of certain crimes.
    \item treating criminals humanely.
    \item the death penalty for the most serious crimes.
\end{enumerate}
\end{enumerate}

\section*{True / False (10 marks)}
\textit{Answer True or False}
\begin{enumerate}[label=\arabic*.]
\setcounter{enumi}{5}
\item True or False: Dershowitz believes that all forms of psychological manipulation techniques are entirely unacceptable in any context.
\item True or False: Anscombe asserts that 'moral obligation' is defined by the laws of a country.
\item True or False: Hasty generalization is a logical fallacy that involves drawing a conclusion based on insufficient evidence.
\item True or False: In Japanese culture, the belief in a merciful God who rewards worshippers is a fundamental theme.
\item True or False: Taurek maintains that the death of two individuals is always worse than the death of one individual.
\end{enumerate}

\section*{Long Form Response (20 marks)}
\textit{Answer each question with a clear and complete explanation.}
\begin{enumerate}[label=\arabic*.]
\setcounter{enumi}{10}
\item Discuss Michael Walzer's four main justifications for terrorism and critically analyze why the notion of terrorism as a form of freedom of speech does not fit within his framework.
\item Discuss the conditions under which a slippery-slope argument is considered fallacious, providing examples to illustrate your analysis.
\end{enumerate}

\vfill
\noindent\textit{End of Paper}
\end{document}